\documentclass[12pt,a4paper]{article}
\usepackage{fullpage}
\usepackage{amsmath,amssymb}
\usepackage{graphicx,xcolor,titlesec,pdfpages}
\usepackage[hidelinks]{hyperref}

\usepackage{natbib}
\bibpunct{(}{)}{;}{a}{}{,}
\setlength{\bibsep}{2pt}
\setlength{\bibhang}{50pt}

% \usepackage{aas_macros}

% \usepackage{bibunits}
\newcommand{\aap}{A\&A}
\newcommand{\aaps}{A\&A Supp}
\newcommand{\aapr}{A\&A Rev}
\newcommand{\apj}{ApJ}
\newcommand{\apjl}{ApJ Lett}
\newcommand{\apjs}{ApJS}
\newcommand{\solphys}{Sol. Phys.}
\newcommand{\etal}{\emph{et al}}
\newenvironment{itemize*}
{\begin{itemize} \setlength{\itemsep}{0pt} \setlength{\parskip}{0pt}}
{\end{itemize}}


\usepackage{soul}
\usepackage{fancyhdr}
\pagestyle{fancy} \fancyhf{}
\renewcommand{\headrulewidth}{0pt}
\fancyfoot[R]{DEMReg/Hannah [\thepage]}

% Changing fonts in overleaf
% see https://www.overleaf.com/learn/latex/font_typefaces
\usepackage[T1]{fontenc}
% \usepackage[utf8]{inputenc}
\usepackage{charter}
% \usepackage{tgtermes}
\usepackage{parskip}
\urlstyle{rm}

\title{Mathematics of the DEMReg approach}
\author{Iain Hannah}
\date{June 01, 2026}


\begin{document}

\maketitle\thispagestyle{fancy}

For more detail, see \citet{2012A&A...539A.146H,2013A&A...553A..10H,2004SoPh..225..293K,1992InvPr...8..849H}.

\vspace{-10pt}
\subsection*{Regularization for the DEM problem}

The observed data value $g_i$ for channel or filter $i$ is related to the underlying differential emission measure (DEM),
$\xi(T)=n^2\,\mathrm{d}h/\mathrm{d}T$ [cm$^{-5}$ K$^{-1}$], and the temperature response $K_i(T)$ by
\begin{equation}\label{eqn:org}
  g_i=\int K_i(T)\,\xi(T)\,\mathrm{d}T.
\end{equation}

After discretising temperature into bins $T_j$, this can be written as a system of linear equations,
\begin{equation}\label{eqn:le}
  g_i=\sum_{j=1}^{N} K_{ij}\,\xi_j,
\end{equation}
where $i=1,\ldots,M$ and $j=1,\ldots,N$. Typically $M<N$; for example, with six SDO/AIA channels and 30 temperature bins we would have $M=6$ and $N=30$.

To solve Eqn.~\ref{eqn:org} and recover the DEM, we seek
\begin{equation}\label{eqn:lsq}
  \boldsymbol{\xi}_{\mathrm{LS}}=\arg\min_{\boldsymbol{\xi}}\left\|\mathbf{K}\boldsymbol{\xi}-\mathbf{g}\right\|^2,
\end{equation}
where $\|\mathbf{x}\|^2=\mathbf{x}^T\mathbf{x}=\sum_{i=1}^{M} x_i^2$, the $L_2$ (Euclidean) norm.

Equation~\ref{eqn:org} defines an ill-posed inverse problem, so additional linear constraints are introduced to suppress noise amplification. A common choice is zero-order regularization, which selects the smallest-norm solution from the infinitely many admissible solutions. Instead of solving Eqn.~\ref{eqn:lsq}, we therefore solve

\begin{equation}
  \arg\min_{\boldsymbol{\xi}}\left\|\mathbf{K}\boldsymbol{\xi}-\mathbf{g}\right\|^2
  \qquad \text{subject to} \qquad
  \left\|\mathbf{L}(\boldsymbol{\xi}-\boldsymbol{\xi}_0)\right\|^2 \le \mathrm{const}.
\end{equation}

This can be solved using Lagrange multipliers, giving

\begin{equation}\label{eqn:lc}
  \boldsymbol{\xi}_\lambda=\arg\min_{\boldsymbol{\xi}}\left(\left\|\mathbf{K}\boldsymbol{\xi}-\mathbf{g}\right\|^2 +
  \lambda\left\|\mathbf{L}(\boldsymbol{\xi}-\boldsymbol{\xi}_0)\right\|^2\right),
\end{equation}

where $\mathbf{L}$ is the constraint matrix, $\lambda$ is the regularization parameter, and $\boldsymbol{\xi}_0$ is an optional initial guess.

\vspace{-10pt}
\subsection*{GSVD to solve the regularized problem}

A solution to Eqn.~\ref{eqn:lc} can be found via the GSVD as a function of $\lambda$ \citep{1992InvPr...8..849H}. For $\mathbf{K}\in \mathbb{R}^{M\times N}$ and $\mathbf{L}\in \mathbb{R}^{N\times N}$, the GSVD yields singular values $\alpha_k$ and $\beta_k$ (with $\phi_k=\alpha_k/\beta_k$) and singular vectors $\mathbf{u}_k$, $\mathbf{v}_k$, and $\mathbf{w}_k$ for $k=1,\ldots,N$, satisfying
$\alpha_k^2+\beta_k^2=1$, $\mathbf{U}^T\mathbf{K}\mathbf{W}=\mathrm{diag}(\boldsymbol{\alpha})$, and $\mathbf{V}^T\mathbf{L}\mathbf{W}=\mathrm{diag}(\boldsymbol{\beta})$. Because in our case $M \leq N$, only the first $M$ singular values $\alpha_k$ are non-zero.

The solution is then, as given in Eqn.~18 of \citet{1992InvPr...8..849H},
\begin{equation}\label{eqn:sol}
\boldsymbol{\xi}_\lambda = \sum_{k=1}^{M}\left(f_k\frac{\mathbf{g}\cdot \mathbf{u}_k}{\alpha_k}
 +(1-f_k)\,\mathbf{w}_k^{-1}\boldsymbol{\xi}_0 \right)\mathbf{w}_k,
\end{equation}
where the filter factors $f_k$ are given by
\begin{equation}
f_k=\frac{\phi_k^2}{\phi_k^2 +\lambda}=\frac{\alpha_k^2}{\alpha_k^2+\lambda\beta_k^2},
\qquad
1-f_k=\frac{\lambda\beta_k^2}{\alpha_k^2+\lambda\beta_k^2}.
\end{equation}

% \begin{equation}\label{eqn:sol}
% \boldsymbol{\xi}_\lambda =\sum_{k=1}^ M\frac{\phi_k^2}{\phi_k^2+\lambda}\left(
% \frac{({\bf g}\cdot {\bf u}_k){\bf w}_k}{\alpha_k}+\frac{\lambda \xi_0\beta_k^2}{\alpha_k^2}
% \right),
% \end{equation}

In the code, Eqn.~\ref{eqn:sol} is implemented as
\begin{equation}\label{eqn:solcode}
 \boldsymbol{\xi}_\lambda = \sum_{k=1}^{M}\frac{\left(\alpha_k(\mathbf{g}\cdot \mathbf{u}_k)+\lambda\beta_k^2\mathbf{w}_k^{-1}\boldsymbol{\xi}_0\right)\mathbf{w}_k}{\alpha_k^2 + \lambda\beta_k^2}.
\end{equation}
This is how it is done in the original IDL code\footnote{\href{https://github.com/ianan/demreg/blob/master/idl_org/dem_inv_reg_solution.pro}{idl\_org/dem\_inv\_reg\_solution.pro}} that accompanied \citet{2012A&A...539A.146H}. In the map versions \citep{2013A&A...553A..10H} of the IDL\footnote{\href{https://github.com/ianan/demreg/blob/261a85156c5a91c77cdc55750c58d0f0dfc7be07/idl/demmap_pos.pro}{idl/demmap\_pos.pro}} and Python\footnote{\href{https://github.com/ianan/demreg/blob/261a85156c5a91c77cdc55750c58d0f0dfc7be07/python/demmap_pos.py}{python/demmap\_pos.py}} code, no initial guess is used and $\boldsymbol{\xi}_\lambda$ is calculated slightly differently, with the regularized inverse of $\mathbf{K}$, $\mathbf{R}_\lambda \simeq \mathbf{K}^{\dagger}$, computed directly as
\begin{equation}\label{eqn:solmapcode}
 \mathbf{R}_\lambda = \sum_{k=1}^{M}\frac{\alpha_k\left(\mathbf{u}_k\mathbf{w}_k\right)}{\alpha_k^2 + \lambda\beta_k^2}.
\end{equation}
The DEM solution is then found from $\boldsymbol{\xi}_\lambda = \mathbf{R}_\lambda\mathbf{g}$. This form is used because $\mathbf{R}_\lambda$ is also needed to estimate the temperature resolution\footnote{\href{https://github.com/ianan/demreg/blob/master/idl_org/dem_inv_reg_resolution.pro}{idl\_org/dem\_inv\_reg\_resolution.pro}} of the DEM solution, so it does not need to be computed separately in the map versions, where computational speed is more important.

Note that:
% \vspace{-10pt}
\begin{itemize*}
\vspace{-5pt}
    \item In Eqn.~6 of \cite{2012A&A...539A.146H}, the $\boldsymbol{\xi}_0$ term is not quite correct, although the code is.
    \item In Eqn.~24 of \cite{2004SoPh..225..293K}, no initial guess is shown, although one is included in the code\footnote{\href{https://hesperia.gsfc.nasa.gov/ssw/packages/xray/idl/inversion/inv_reg_solution.pro}{ssw/packages/xray/idl/inversion/inv\_reg\_solution.pro}}.
    \item In Eqn.~15 of \citet{1992InvPr...8..849H}, the regularization parameter is written as $\lambda^2$ rather than $\lambda$, as used here.
\end{itemize*} 

  To obtain a useful solution, we still need to choose suitable values of $\mathbf{L}$ and $\lambda$.

%\newpage
%\vspace{-10pt}
\subsection*{Choosing constraint matrix ${\bf L}$}

In most applications we consider zero-order regularization, for which $\mathbf{L} \propto \mathbf{I}$ and therefore has no off-diagonal terms. The diagonal elements of $\mathbf{L}$ are usually normalized by some scale, typically either the temperature-bin width, $\mathbf{L}_{jj}=1/\mathrm{d}\log T_j$, or, when appropriate, an expected form for the final solution (which is not the same as the initial guess $\boldsymbol{\xi}_0$). One example is the minimum of the EM loci curves, $\boldsymbol{\xi}_{\mathrm{mlc}}$, taken as a smoothed and normalized version of the lowest emission measure at each temperature under an isothermal assumption; in that case $\mathbf{L}_{jj}$ is based on the minimum of $g_i/K_{ij}$ over all $i$.


\vspace{-10pt}
\subsection*{Choosing regularization parameter $\lambda$}

The regularization parameter $\lambda$ is chosen so that the solution $\boldsymbol{\xi}_\lambda$ from Eqn.~\ref{eqn:sol} satisfies an additional criterion. The standard choice is Morozov's discrepancy principle, which effectively controls the $\chi^2$ of the solution in data space. If the observations have associated uncertainties $\delta g_i$, then the preferred value of $\lambda$ is the one that minimizes
\begin{equation}\label{eqn:desc}
  \left\|\mathbf{K}\boldsymbol{\xi}_\lambda-\mathbf{g}\right\|^2 - \rho\left\|\boldsymbol{\delta g}\right\|^2.
\end{equation}

where $\rho$ is a regularization-tuning parameter; by default we aim for $\rho=1$.



%\begin{equation}
%    \left\|\bf{K\boldsymbol{\xi}_\lambda-g}\right\|^2  =  \rho \left\| {\bf \delta g}\right\|^2\;,
%\end{equation}
%
%
%
%
%\begin{equation}
%    \frac{1}{M}\left\|\bf{K\boldsymbol{\xi}_\lambda-g}\right\|^2  =  \rho
%\end{equation}
%where $\rho$ is the regularization tweak parameter, which is effectively controlling the $\chi^2$ of the solution in data space. So by default are aiming for $\rho=1$. 
%So the optimal $\lambda$ should satisfy 
%\begin{equation}\label{eqn:desc}
%    \left\|\bf{K\boldsymbol{\xi}_\lambda-g}\right\|^2  -  \rho \left\| {\bf \delta g}\right\|^2 = \mathrm{min}.
%\end{equation}

From Eqn.~20 of \citet{1992InvPr...8..849H}, the term $\left\|\mathbf{K}\boldsymbol{\xi}_\lambda-\mathbf{g}\right\|^2$ can be calculated directly from the GSVD products, namely
\begin{equation}\label{eqn:sol2}
\left\|\mathbf{K}\boldsymbol{\xi}_\lambda-\mathbf{g}\right\|^2=\sum_{k=1}^{M} (1-f_k)^2\left(
\mathbf{g}\cdot \mathbf{u}_k - \alpha_k\mathbf{w}_k^{-1}\boldsymbol{\xi}_0
\right)^2.
\end{equation}

In the code, Eqn.~\ref{eqn:sol2} is implemented as
\begin{equation}\label{eqn:finarg}
  \left\|\mathbf{K}\boldsymbol{\xi}_\lambda-\mathbf{g}\right\|^2=\sum_{k=1}^{M} \left(\frac{\lambda\beta_k^2\left(\mathbf{g}\cdot \mathbf{u}_k-\alpha_k\mathbf{w}_k^{-1}\boldsymbol{\xi}_0\right)}{\alpha_k^2 + \lambda\beta_k^2}\right)^2.
\end{equation}

This is how it is done in the original IDL code\footnote{\href{https://github.com/ianan/demreg/blob/master/idl_org/dem_inv_reg_parameter.pro}{idl\_org/dem\_inv\_reg\_parameter.pro}}. In the map versions of the IDL\footnote{\href{https://github.com/ianan/demreg/blob/master/idl/dem_inv_reg_parameter_map.pro}{idl/dem\_inv\_reg\_parameter\_map.pro}} and Python\footnote{\href{https://github.com/ianan/demreg/blob/master/python/dem_reg_map.py}{python/dem\_reg\_map.py}} code, Eqn.~\ref{eqn:finarg} should have used only $\left(\mathbf{g}\cdot \mathbf{u}_k\right)$ in the numerator because no initial guess, i.e. no $\boldsymbol{\xi}_0$, is used. However, $\left(\mathbf{g}\cdot \mathbf{u}_k-\alpha_k\right)$ was used instead, which may have affected the determination of the optimal $\lambda$. Crucially, the calculation of $\boldsymbol{\xi}_\lambda$ itself was correct in all versions. This has now been corrected in the map versions of the code.

In all versions of the code, Eqn.~\ref{eqn:finarg} is evaluated over a range of $\lambda$ values, with the range set by scaling the minimum and maximum of the $\phi_k$ values. The left-hand side of Eqn.~\ref{eqn:desc} is then computed for each case, and the $\lambda$ giving the smallest value is selected.

Using Eqn.~\ref{eqn:desc} should provide the mathematically optimal solution, but not necessarily the most physical one, because it can return a $\boldsymbol{\xi}_\lambda$ with negative components. An additional positivity constraint can therefore be imposed, requiring the chosen $\lambda$ to satisfy Eqn.~\ref{eqn:desc} while also giving $\boldsymbol{\xi}_\lambda \ge 0$ component-wise. In the original IDL code\footnote{\href{https://github.com/ianan/demreg/blob/master/idl_org/dem_inv_reg_parameter_pos.pro}{idl\_org/dem\_inv\_reg\_parameter\_pos.pro}}, this is achieved by evaluating both Eqn.~\ref{eqn:finarg} and Eqn.~\ref{eqn:solcode} for a range of $\lambda$ values and selecting the one that satisfies both criteria. In the map code (IDL\footnote{\href{https://github.com/ianan/demreg/blob/master/idl/demmap_pos.pro}{idl/demmap\_pos.pro}} and Python\footnote{\href{https://github.com/ianan/demreg/blob/master/python/demmap_pos.py}{python/demmap\_pos.py}}), Eqn.~\ref{eqn:finarg} was evaluated for a smaller number of $\lambda$ samples and, if $\boldsymbol{\xi}_\lambda$ contained negative components, $\rho$ was increased iteratively and Eqn.~\ref{eqn:finarg} recalculated until $\boldsymbol{\xi}_\lambda \ge 0$ or a maximum number of iterations was reached. This was faster than evaluating both Eqn.~\ref{eqn:finarg} and Eqn.~\ref{eqn:solcode} over a much larger $\lambda$ sample, as in the original approach. A possible optimization of the original method would be to first identify a subset of $\lambda$ values satisfying one criterion and then apply the second criterion only to that smaller subset; for example, first find the values that give $\boldsymbol{\xi}_\lambda \ge 0$ and then evaluate Eqn.~\ref{eqn:finarg} only on that subset.

\subsection*{Errors on the DEM solution}

For more details on how these are calculated see \citet{2012A&A...539A.146H,2013A&A...553A..10H} and the code. But in summary:
\begin{itemize}
	\item The vertical errors on the DEM solution are obtained from linear propagation of the uncertainties in the input data.
	\item The horizontal errors on the DEM solution correspond to the temperature resolution and are estimated from the spread of the off-diagonal terms in $\mathbf{K}\mathbf{R}_\lambda$. The better the solution $\boldsymbol{\xi}_\lambda = \mathbf{R}_\lambda\mathbf{g}$, the closer $\mathbf{K}\mathbf{R}_\lambda$ is to $\mathbf{I}$.
\end{itemize}

\subsection*{Temperature response functions}

The temperature response functions are given by
\begin{equation}
  K_i(T_j)=\int_0^\infty G(\lambda,T_j)\,R_i(\lambda)\,\mathrm{d}\lambda,
\end{equation}
where $G(\lambda,T)$ is the plasma emissivity in the solar atmosphere (the contribution function), derived from an appropriate solar model such as CHIANTI, and $R_i(\lambda)$ is the wavelength response of the $i^{\mathrm{th}}$ channel. These temperature response functions are usually precomputed and provided by the instrument teams, for example for SDO/AIA \citep{2014SoPh..289.2377B}. With calibrated EUV spectral-line data, however, the contribution functions $G_i(T_j)$ are used directly and can be calculated with CHIANTI's gofnt\footnote{i.e. \href{https://hesperia.gsfc.nasa.gov/ssw/packages/chianti/idl/emiss/gofnt.pro}{SSW version gofnt.pro}} function.

\newpage
\par\noindent\rule{0.5\textwidth}{0.4pt}

%\bibliographystyle{aasjournal}
%\bibliography{refs.bib}	
\begin{thebibliography}{}
\expandafter\ifx\csname natexlab\endcsname\relax\def\natexlab#1{#1}\fi
\providecommand{\url}[1]{\href{#1}{#1}}
\providecommand{\dodoi}[1]{doi:~\href{http://doi.org/#1}{\nolinkurl{#1}}}
\providecommand{\doeprint}[1]{\href{http://ascl.net/#1}{\nolinkurl{http://ascl.net/#1}}}
\providecommand{\doarXiv}[1]{\href{https://arxiv.org/abs/#1}{\nolinkurl{https://arxiv.org/abs/#1}}}

\bibitem[{{Boerner} {et~al.}(2014){Boerner}, {Testa}, {Warren}, {Weber}, \&
  {Schrijver}}]{2014SoPh..289.2377B}
{Boerner}, P.~F., {Testa}, P., {Warren}, H., {Weber}, M.~A., \& {Schrijver},
  C.~J. 2014, \solphys, 289, 2377, \dodoi{10.1007/s11207-013-0452-z}

\bibitem[{{Hannah} \& {Kontar}(2012)}]{2012A&A...539A.146H}
{Hannah}, I.~G., \& {Kontar}, E.~P. 2012, \aap, 539, A146,
  \dodoi{10.1051/0004-6361/201117576}

\bibitem[{{Hannah} \& {Kontar}(2013)}]{2013A&A...553A..10H}
---. 2013, \aap, 553, A10, \dodoi{10.1051/0004-6361/201219727}

\bibitem[{{Hansen}(1992)}]{1992InvPr...8..849H}
{Hansen}, P.~C. 1992, Inverse Problems, 8, 849,
  \dodoi{10.1088/0266-5611/8/6/005}

\bibitem[{{Kontar} {et~al.}(2004){Kontar}, {Piana}, {Massone}, {Emslie}, \&
  {Brown}}]{2004SoPh..225..293K}
{Kontar}, E.~P., {Piana}, M., {Massone}, A.~M., {Emslie}, A.~G., \& {Brown},
  J.~C. 2004, \solphys, 225, 293, \dodoi{10.1007/s11207-004-4140-x}

\end{thebibliography}


\end{document}
